\chapter{Formato dell'esame: cosa sappiamo}
\label{cap:formato}

\section{Struttura (da fonte non ufficiale)}

L'unica traccia disponibile è un messaggio nel gruppo Telegram del corso, presentato come elenco
delle domande di un appello precedente (Capitolo~\ref{cap:tema-precedente}). Da quel messaggio
risulta questa struttura:
\begin{itemize}
  \item \textbf{4 domande chiuse} a risposta multipla, alcune nella forma ``\emph{select the WRONG
    answer}'' (individuare l'affermazione errata);
  \item \textbf{1 domanda aperta} su un caso pratico, divisa in sotto-domande: quale norma si applica,
    cosa va fatto secondo la norma, quali obblighi valgono a prescindere dall'incidente;
  \item domande in \textbf{inglese}.
\end{itemize}

\begin{verifica}
\begin{itemize}
  \item È \textbf{una sola fonte}, non verificata. Per la stessa regola dell'audit spiegata in aula,
    una fonte sola è un'asserzione, non un fatto: va confermata con il docente o con altri studenti.
  \item L'appello si riferisce con ogni probabilità a un anno accademico precedente. Allora il CRA non
    era ancora operativo (obblighi di segnalazione dall'11/9/2026), mentre quest'anno è stato centrale:
    è ragionevole aspettarsi domande sul CRA e su casi CRA--GDPR.
\end{itemize}
\end{verifica}

\section{Indicazioni date dai docenti durante il corso}

\begin{itemize}
  \item L'esercizio di \textbf{notifica alle autorità} è ``molto probabile'' (lezione del 02/11).
  \item Si valuta la capacità di \textbf{qualificare il caso} (quale norma) e di indicare il
    \textbf{percorso corretto} in caso di incidente; non serve ricordare a memoria gli allegati.
  \item Le \textbf{scadenze} vanno ricordate con precisione, senza confondere le fasi delle diverse
    norme (Capitolo~\ref{cap:scadenze}). Il professore ha parlato di una tabella riassuntiva
    ``prodotta per voi per l'esame''.
  \item I \textbf{dettagli del caso} contano (es.\ una DPIA mai aggiornata): vanno letti e usati nella
    risposta.
  \item Nelle domande aperte si premia il \textbf{ragionamento motivato} (necessità, proporzionalità,
    accountability), non l'elenco di articoli.
\end{itemize}

\begin{verifica}
Non è chiaro se la tabella riassuntiva delle scadenze sarà consultabile durante l'esame o se è solo
materiale di studio. Chiedere al docente.
\end{verifica}

\section{Metodo per la domanda aperta}

\begin{enumerate}
  \item \textbf{Qualifica}: elenca le norme applicabili \emph{e} spiega in una riga perché le altre
    non si applicano (es.\ ``DORA no: non è un'entità finanziaria'').
  \item \textbf{Ruoli}: chi è titolare, responsabile, soggetto essenziale/importante, fabbricante.
  \item \textbf{Orologi}: per ogni norma, da quando decorre e verso quale autorità.
  \item \textbf{Severità}: valuta il rischio per gli interessati (ENISA) e decidi su art.~33 e 34.
  \item \textbf{Misure}: contenimento, comunicazioni, registri, misure organizzative e contrattuali.
  \item \textbf{Usa i dettagli del caso} per motivare ogni scelta.
\end{enumerate}
